% Lecture notes: Week 10, Lecture 3 -- Stress Transformations with Mohr's Circle (asynchronous)
% To hide this lecture from the website, add a line containing only:  % draft
% To set the date shown on the website, add a line like:  % posted: 2026-09-02
\documentclass[11pt]{article}
\usepackage{mech2030notes}
\begin{document}

% ------------------------------------------------------------------
% Lecture-local helpers (defined inside the document so they also work
% in the combined notes; names are unique to this lecture).
%
% \WTenLThreeMohr{n}: Mohr's circle for Problem 1 after stage n, with the
% piece added at stage n drawn in red.
% Scale: 0.3 cm per ksi; tau is plotted positive DOWNWARD.
%   A  = (sigma_x, tau_xy)   = (6, -4)    -> (1.8, 1.2)
%   G  = (sigma_y, -tau_xy)  = (-10, 4)   -> (-3, -1.2)
%   C  = (sigma_avg, 0)      = (-2, 0)    -> (-0.6, 0),  R = 2.6833
%   X' = (sigma_x', tau_x'y'): NOT to scale. The true 2*theta = 150 deg would put X'
%   only 3.4 deg above the sigma-axis; it is drawn at 2*theta ~ 135 deg (screen angle
%   162 deg) so that the small angle alpha (~18 deg as drawn) is visible.
%        -> X' = (-3.152, 0.829),  Y' = (1.952, -0.829)
% ------------------------------------------------------------------
\newcommand{\WTenLThreeCol}[2]{\ifnum#1=#2\def\wtlc{red!75!black}\else\def\wtlc{black}\fi}
\newcommand{\WTenLThreeMohr}[1]{%
  \ifnum#1>1
    \WTenLThreeCol{#1}{2}%
    \draw[axis, \wtlc] (-4.5,0) -- (3.5,0) node[right] {$\sigma$};
    \draw[axis, \wtlc] (0,3.1) -- (0,-3.3) node[below] {$\tau$};
    \foreach \s in {-3.6,-2.4,-1.2,1.2,2.4} \draw[thin, \wtlc] (\s,0.07) -- (\s,-0.07);
  \fi
  \ifnum#1>5
    \WTenLThreeCol{#1}{6}%
    \draw[semithick, \wtlc] (-0.6,0) circle (2.6833);
  \fi
  \ifnum#1>4
    \WTenLThreeCol{#1}{5}%
    \draw[thick, \wtlc] (-3,-1.2) -- (1.8,1.2);
    \node[pin, \wtlc] at (-0.6,0) {};
    \node[below right, \wtlc] at (-0.6,0) {$\sigma_{\text{avg}}$};
  \fi
  \ifnum#1>2
    \WTenLThreeCol{#1}{3}%
    \node[pin, \wtlc] at (1.8,1.2) {};
    \node[above right, \wtlc] at (1.8,1.2) {$A\,(\sigma_x, \tau_{xy})$};
  \fi
  \ifnum#1>3
    \WTenLThreeCol{#1}{4}%
    \node[pin, \wtlc] at (-3,-1.2) {};
    \node[below left, \wtlc] at (-3,-1.2) {$G\,(\sigma_y, -\tau_{xy})$};
  \fi
  \ifnum#1>6
    \WTenLThreeCol{#1}{7}%
    \draw[thick, \wtlc] (-0.6,0) -- (-0.6,-2.6833);
    \node[right, \wtlc] at (-0.6,-1.9) {$R = \tau_{\max}$};
  \fi
  \ifnum#1>7
    \WTenLThreeCol{#1}{8}%
    \draw[thinarrow, \wtlc] (-0.6,0) ++(26.565:1.15) arc (26.565:162:1.15);
    \node[\wtlc] at ($(-0.6,0)+(100:1.5)$) {$2\theta$};
    \draw[very thick, \wtlc] (-3.152,0.829) -- (1.952,-0.829);
    \node[pin, \wtlc] at (-3.152,0.829) {};
    \node[above left, \wtlc] at (-3.152,0.829) {$X'(\sigma_{x'}, \tau_{x'y'})$};
  \fi
  \ifnum#1>8
    \WTenLThreeCol{#1}{9}%
    \draw[thinarrow, \wtlc] (-0.6,0) ++(26.565:0.75) arc (26.565:0:0.75);
    \node[\wtlc] at ($(-0.6,0)+(12:1.25)$) {$2\theta_{p1}$};
    \draw[densely dotted, \wtlc] (-3.152,0.829) -- (-3.152,0);
    \node[pin, \wtlc] at (-3.2833,0) {};
    \draw[thick, \wtlc] (-0.6,0) ++(162:1.75) arc (162:180:1.75);
    \node[\wtlc] at ($(-0.6,0)+(171:2.1)$) {$\alpha$};
  \fi
}
\newcommand{\WTenLThreePanel}[1]{%
  \begin{tikzpicture}[scale=0.85, every node/.style={font=\scriptsize}]
    \WTenLThreeMohr{#1}
  \end{tikzpicture}}

% ##################################################################
%  WEEK 10, LECTURE 3
% ##################################################################
\lecture{10}{3}{Stress Transformations with Mohr's Circle}

\begin{center}
\fbox{\parbox{0.9\textwidth}{\centering
  \textbf{Asynchronous session.} Work through Problem~1 below step by step, then try
  Problems~2 and~3 on your own and check your answers.}}
\end{center}

\noindent
$\rightarrow$ \textbf{Question for today:} how do we use Mohr's circle to compute the
stresses acting on a \emph{particular} plane, cut at a given angle through the material?

\subsection*{What you will need (from Week~10, Lectures~1--2)}

\noindent
\begin{minipage}[t]{0.5\textwidth}
\begin{align*}
  \sigma_{\text{avg}} &= \frac{\sigma_x + \sigma_y}{2} \\
  R = \tau_{\max} &= \sqrt{\left(\frac{\sigma_x - \sigma_y}{2}\right)^2 + \tau_{xy}^2} \\
  \theta_{p1} &= \frac12 \tan^{-1}\!\left(\frac{2\tau_{xy}}{\sigma_x - \sigma_y}\right)
\end{align*}
\end{minipage}%
\hfill
\begin{minipage}[t]{0.46\textwidth}
\textbf{Key rule:} rotating the element by an angle $\theta$
(counterclockwise $= $ positive) corresponds to rotating by $2\theta$ \emph{in the same
direction} on Mohr's circle (with $\tau$ plotted positive \emph{downward}).

\medskip
The point $A = (\sigma_x, \tau_{xy})$ represents the $x$-face, i.e.\ $\theta = \ang{0}$.
\end{minipage}

% ==================================================================
\section*{Problem 1}
% ==================================================================

\noindent
\begin{minipage}[c]{0.55\textwidth}
\textbf{Setup:} The stress state shown acts on a piece of solid material.

\smallskip
\textbf{Find:} The normal stress $\sigma_{x'}$ and shear stress $\tau_{x'y'}$ on the
inclined (cut) face shown.

\smallskip
\textbf{Steps:}
\begin{enumerate}[itemsep=1pt]
  \item Draw Mohr's circle.
  \item What angle $\theta$ does the rotation shown correspond to?
  \item Draw a diameter oriented at $2\theta$ on your Mohr's circle.
  \item Use that diameter to determine $\sigma_{x'}$ and $\tau_{x'y'}$.
\end{enumerate}
\end{minipage}%
\hfill
\begin{minipage}[c]{0.42\textwidth}
\centering
\begin{tikzpicture}
  \draw[thick, fill=gray!12] (0,0) rectangle (2,2);
  \draw[thick] (0,2) -- (2,1.4641);
  \draw[thinarrow] (0,1.45) arc (270:345:0.55);
  \node at (0.45,1.3) {\small$\ang{75}$};
  % sigma_x = +6 ksi
  \draw[force] (2.05,1) -- (2.8,1) node[above right] {$\SI{6}{ksi}$};
  \draw[force] (-0.05,1) -- (-0.8,1);
  % sigma_y = -10 ksi (compression)
  \draw[force] (1,2.8) node[above] {$\SI{10}{ksi}$} -- (1,2.05);
  \draw[force] (1,-0.8) -- (1,-0.05);
  % tau_xy = -4 ksi
  \draw[force] (2.25,1.6) -- (2.25,0.4) node[below right] {$\SI{4}{ksi}$};
  \draw[force] (-0.25,0.4) -- (-0.25,1.6);
  \draw[force] (1.6,2.25) -- (0.4,2.25);
  \draw[force] (0.4,-0.25) -- (1.6,-0.25);
\end{tikzpicture}
\end{minipage}

% ------------------------------------------------------------------
\subsection*{Step 0: What are $\sigma_{x'}$ and $\tau_{x'y'}$?}
% ------------------------------------------------------------------

\noindent
Before drawing anything, let's be clear about what we are looking for. Imagine cutting the
element along the inclined line and keeping the lower piece. The stresses we want act on
the newly exposed cut face:

\noindent
\begin{minipage}[c]{0.5\textwidth}
\centering
\begin{tikzpicture}[scale=1.45]
  % lower piece; cut face midpoint M = (1, 1.732); outward normal n at 75 deg
  \coordinate (M) at (1,1.732);
  \draw[thick, fill=gray!12] (0,0) -- (2,0) -- (2,1.4641) -- (0,2) -- cycle;
  % stresses on the original faces
  \draw[force] (2.05,0.73) -- (2.7,0.73) node[right] {\small$\SI{6}{ksi}$};
  \draw[force] (-0.05,1) -- (-0.7,1);
  \draw[force] (1,-0.7) node[below] {\small$\SI{10}{ksi}$} -- (1,-0.05);
  \draw[force] (2.22,1.25) -- (2.22,0.3) node[below right] {\small$\SI{4}{ksi}$};
  \draw[force] (-0.22,0.4) -- (-0.22,1.6);
  \draw[force] (0.4,-0.22) -- (1.6,-0.22);
  % local axes (dashed) along the normal and along the face
  \draw[thin, dashed] ($(M)+(75:1.05)$) -- ($(M)+(75:1.5)$) node[above] {$x'$};
  \draw[thin, dashed] ($(M)+(165:0.75)$) -- ($(M)+(165:1.4)$) node[left] {$y'$};
  % unknowns on the cut face (drawn in their positive directions)
  \draw[force, red!75!black] ($(M)+(75:0.05)$) -- ($(M)+(75:1.0)$);
  \node[red!75!black, right] at ($(M)+(75:0.7)$) {$\sigma_{x'}$};
  \draw[force, red!75!black] ($(M)+(75:0.22)+(165:-0.1)$) -- ($(M)+(75:0.22)+(165:0.7)$);
  \node[red!75!black, above] at ($(M)+(75:0.3)+(165:0.55)$) {$\tau_{x'y'}$};
\end{tikzpicture}
\end{minipage}%
\hfill
\begin{minipage}[c]{0.47\textwidth}
\begin{itemize}
  \item The $x'$ axis is the \emph{outward normal} to the cut face.
  \item The $y'$ axis lies \emph{along} the cut face ($x'$ rotated $\ang{90}$
        counterclockwise).
  \item $\sigma_{x'}$ acts along $x'$, perpendicular to the face. It is positive in
        \emph{tension} (pulling away from the face).
  \item $\tau_{x'y'}$ acts along the face. It is positive when it points in the $+y'$
        direction (on a face whose outward normal is $+x'$).
\end{itemize}
The red arrows show both unknowns in their \emph{positive} directions. If an answer comes
out negative, the stress actually points the other way.
\end{minipage}

% ------------------------------------------------------------------
\subsection*{Step 1: Draw Mohr's circle}
% ------------------------------------------------------------------

\noindent
We build the circle in small pieces. Each panel shows the circle after that piece; the piece
just added is in \textcolor{red!75!black}{red}. (Scale: one tick $= \SI{4}{ksi}$.)

\medskip
\noindent
\textbf{1a. Read off the stress state} from the element (watch the signs!):
\begin{itemize}
  \item $\sigma_x = +\SI{6}{ksi}$: the arrows on the $x$-faces point \emph{away} from the
        element (tension).
  \item $\sigma_y = -\SI{10}{ksi}$: the arrows on the $y$-faces point \emph{into} the
        element (compression).
  \item $\tau_{xy} = -\SI{4}{ksi}$: on the right ($+x$) face the shear points \emph{down}
        ($-y$), and on the top ($+y$) face it points \emph{left} ($-x$). Both are opposite
        to the positive convention, so $\tau_{xy}$ is negative.
\end{itemize}

\bigskip
\noindent
\begin{minipage}[c]{0.45\textwidth}
\textbf{1b. Draw the axes.}
\begin{itemize}
  \item Horizontal: $\sigma$, positive to the right.
  \item Vertical: $\tau$, positive \emph{downward}.
\end{itemize}
\end{minipage}%
\hfill
\begin{minipage}[c]{0.52\textwidth}
\centering\WTenLThreePanel{2}
\end{minipage}

\bigskip
\noindent
\begin{minipage}[c]{0.45\textwidth}
\textbf{1c. Draw a dot at $(\sigma_x, \tau_{xy})$:}
\[ A = (6,\ -4)\ \si{ksi} \]
Because $\tau_{xy} < 0$ and $\tau$ is plotted positive downward, point $A$ lies
\emph{above} the $\sigma$-axis.
\end{minipage}%
\hfill
\begin{minipage}[c]{0.52\textwidth}
\centering\WTenLThreePanel{3}
\end{minipage}

\bigskip
\noindent
\begin{minipage}[c]{0.45\textwidth}
\textbf{1d. Draw a dot at $(\sigma_y, -\tau_{xy})$:}
\[ G = (-10,\ 4)\ \si{ksi} \]
This one lies \emph{below} the $\sigma$-axis.
\end{minipage}%
\hfill
\begin{minipage}[c]{0.52\textwidth}
\centering\WTenLThreePanel{4}
\end{minipage}

\bigskip
\noindent
\begin{minipage}[c]{0.45\textwidth}
\textbf{1e. Draw a line between the two dots.} Where it crosses the $\sigma$-axis is
$\sigma_{\text{avg}}$:
\[ \sigma_{\text{avg}} = \frac{\sigma_x + \sigma_y}{2} = \frac{6 + (-10)}{2} \]
\[ \boxed{\sigma_{\text{avg}} = -\SI{2}{ksi}} \]
\end{minipage}%
\hfill
\begin{minipage}[c]{0.52\textwidth}
\centering\WTenLThreePanel{5}
\end{minipage}

\bigskip
\noindent
\begin{minipage}[c]{0.45\textwidth}
\textbf{1f. Draw the circle} centered at $(\sigma_{\text{avg}}, 0)$ and passing through
$A$ and $G$.
\begin{itemize}[label=\then]
  \item Every point on this circle is the stress state on some rotated face.
\end{itemize}
\end{minipage}%
\hfill
\begin{minipage}[c]{0.52\textwidth}
\centering\WTenLThreePanel{6}
\end{minipage}

\bigskip
\noindent
\begin{minipage}[c]{0.45\textwidth}
\textbf{1g. What is the radius of the circle?} This radius is $\tau_{\max}$:
\begin{align*}
  R &= \sqrt{\left(\frac{6 - (-10)}{2}\right)^2 + (-4)^2} \\
    &= \sqrt{8^2 + 4^2}
\end{align*}
\[ \boxed{R = \tau_{\max} = \SI{8.94}{ksi}} \]
\end{minipage}%
\hfill
\begin{minipage}[c]{0.52\textwidth}
\centering\WTenLThreePanel{7}
\end{minipage}

% ------------------------------------------------------------------
\subsection*{Step 2: What angle $\theta$ does the rotation shown correspond to?}
% ------------------------------------------------------------------

\noindent
\begin{minipage}[c]{0.45\textwidth}
\centering
\begin{tikzpicture}[scale=1.3]
  \draw[thick, fill=gray!12] (0,0) rectangle (2,2);
  \draw[thick] (0,2) -- (2,1.4641);
  \draw[thinarrow] (0,1.5) arc (270:345:0.5);
  \node at (0.42,1.33) {\scriptsize$\ang{75}$};
  % local axes at a point P on the cut
  \coordinate (P) at (1,1.732);
  \draw[axis] (P) -- ++(1.35,0) node[right] {$x$};
  \draw[axis, red!75!black] (P) -- ++(75:1.25) node[above] {$x'$};
  \draw[thinarrow, red!75!black] ($(P)+(0.75,0)$) arc (0:75:0.75);
  \node[red!75!black] at ($(P)+(38:1.05)$) {$\theta$};
  \draw ($(P)+(1.05,0)$) arc (0:-15:1.05);
  \node at ($(P)+(-8:1.35)$) {\scriptsize$\ang{15}$};
\end{tikzpicture}
\end{minipage}%
\hfill
\begin{minipage}[c]{0.52\textwidth}
\begin{itemize}
  \item The cut makes $\ang{75}$ with the vertical left edge, so it sits
        $\ang{90} - \ang{75} = \ang{15}$ \emph{below} the horizontal.
  \item The outward normal $x'$ is perpendicular to the cut, so it sits
        $\ang{90} - \ang{15} = \ang{75}$ \emph{above} the $x$ axis.
  \item Going from $x$ to $x'$ is a \emph{counterclockwise} rotation, so $\theta$ is
        positive:
\end{itemize}
\[
  \boxed{\theta = +\ang{75}}
  \qquad\Rightarrow\qquad
  \boxed{2\theta = \ang{150}}
\]
\end{minipage}

% ------------------------------------------------------------------
\subsection*{Step 3: Draw a diameter oriented at $2\theta$ on your Mohr's circle}
% ------------------------------------------------------------------

\noindent
\begin{minipage}[c]{0.45\textwidth}
\begin{itemize}
  \item Start from point $A$ (the $x$-face, $\theta = \ang{0}$).
  \item Rotate by $2\theta = \ang{150}$ around the center, \emph{counterclockwise}
        (the same direction as the element rotation).
  \item The point you land on is $X' = (\sigma_{x'}, \tau_{x'y'})$, the stress state on
        the cut face. Draw the new diameter through $X'$ and the center.
\end{itemize}
\end{minipage}%
\hfill
\begin{minipage}[c]{0.52\textwidth}
\centering\WTenLThreePanel{8}
\end{minipage}

% ------------------------------------------------------------------
\subsection*{Step 4: Use that diameter to determine $\sigma_{x'}$ and $\tau_{x'y'}$}
% ------------------------------------------------------------------

\noindent
\begin{minipage}[c]{0.45\textwidth}
\textbf{4a. Find where the original diameter sits.} The angle from $A$ down to the
$\sigma$-axis is $2\theta_{p1}$:
\[
  \theta_{p1} = \frac12\tan^{-1}\!\left(\frac{2(-4)}{6-(-10)}\right) = -\ang{13.28}
\]
so $|2\theta_{p1}| = \ang{26.57}$.
\end{minipage}%
\hfill
\begin{minipage}[c]{0.52\textwidth}
\centering\WTenLThreePanel{9}
\end{minipage}

\bigskip
\noindent
\textbf{4b. Find the small angle $\alpha$} between the new diameter and the $\sigma$-axis.
Going all the way around the top of the circle from $A$ to the left end of the
$\sigma$-axis is $\ang{180}$, made up of $|2\theta_{p1}|$, then $2\theta$, then $\alpha$:
\[
  \alpha = \ang{180} - |2\theta| - |2\theta_{p1}| = \ang{180} - \ang{150} - \ang{26.57}
  \qquad\Rightarrow\qquad
  \boxed{\alpha = \ang{3.43}}
\]

\noindent
\textbf{4c. Read $\sigma_{x'}$ and $\tau_{x'y'}$ off the triangle} formed by the center,
$X'$, and the $\sigma$-axis:

\noindent
\begin{minipage}[c]{0.45\textwidth}
\centering
\begin{tikzpicture}[scale=1.1]
  % not to scale: alpha exaggerated for clarity
  \draw[axis] (-0.6,0) -- (4.0,0) node[right] {$\sigma$};
  \draw[axis] (3.4,0.9) -- (3.4,-0.7) node[below] {$\tau$};
  \coordinate (C) at (3,0);
  \coordinate (X) at (0,0.8);
  \draw[very thick, red!75!black] (C) -- (X) node[midway, above, black] {$R$};
  \node[pin] at (C) {}; \node[below] at (C) {$\sigma_{\text{avg}}$};
  \node[pin, red!75!black] at (X) {}; \node[above] at (X) {$X'$};
  \draw[densely dotted] (X) -- (0,0);
  \draw (2.2,0) arc (180:165.07:0.8);
  \node at (1.95,0.11) {\small$\alpha$};
  \draw[dim] (0,-0.45) -- (3,-0.45) node[midway, below] {$R\cos\alpha$};
  \draw[dim] (-0.35,0) -- (-0.35,0.8) node[midway, left] {$R\sin\alpha$};
\end{tikzpicture}\\[-0.2em]
{\scriptsize (not to scale: $\alpha$ is exaggerated)}
\end{minipage}%
\hfill
\begin{minipage}[c]{0.52\textwidth}
\begin{itemize}
  \item $X'$ is a distance $R\cos\alpha$ to the \emph{left} of the center:
        \[
          \sigma_{x'} = \sigma_{\text{avg}} - R\cos\alpha
          = -2 - 8.94\cos(\ang{3.43})
        \]
  \item $X'$ is a distance $R\sin\alpha$ \emph{above} the $\sigma$-axis. Since $\tau$ is
        plotted positive downward, this is a \emph{negative} shear stress:
        \[
          \tau_{x'y'} = -R\sin\alpha = -8.94\sin(\ang{3.43})
        \]
\end{itemize}
\end{minipage}

\begin{center}
\fbox{\fbox{$\sigma_{x'} = -\SI{10.93}{ksi}, \qquad \tau_{x'y'} = -\SI{0.536}{ksi}$}}
\end{center}

% ------------------------------------------------------------------
\subsection*{What does the answer mean?}
% ------------------------------------------------------------------

\noindent
\begin{minipage}[c]{0.45\textwidth}
\centering
\begin{tikzpicture}[scale=1.45]
  \coordinate (M) at (1,1.732);
  \draw[thick, fill=gray!12] (0,0) -- (2,0) -- (2,1.4641) -- (0,2) -- cycle;
  \draw[thin, dashed] ($(M)+(75:1.3)$) -- ($(M)+(75:1.6)$) node[above] {$x'$};
  \draw[thin, dashed] ($(M)+(165:0.7)$) -- ($(M)+(165:1.3)$) node[left] {$y'$};
  % sigma_x' < 0: compression, the arrow pushes INTO the face
  \draw[force, red!75!black] ($(M)+(75:1.25)$) -- ($(M)+(75:0.05)$);
  \node[red!75!black, right, align=left] at ($(M)+(75:1.0)$) {\small$\SI{10.93}{ksi}$\\[-2pt]\scriptsize(compression)};
  % tau_x'y' < 0: points in the -y' direction (down and to the right along the face)
  \draw[force, red!75!black] ($(M)+(75:0.22)+(165:0.1)$) -- ($(M)+(75:0.22)+(165:-0.65)$);
  \node[red!75!black, right] at ($(M)+(75:0.22)+(165:-0.68)$) {\small$\SI{0.536}{ksi}$};
\end{tikzpicture}
\end{minipage}%
\hfill
\begin{minipage}[c]{0.52\textwidth}
\begin{itemize}
  \item $\sigma_{x'} < 0$: the cut face is in \emph{compression}, so the arrow pushes into
        the face.
  \item $\tau_{x'y'} < 0$: the shear points in the $-y'$ direction, down and to the
        right along the face. It is small.
\end{itemize}
\end{minipage}

\bigskip
\noindent
\textbf{Sanity checks.}
\begin{itemize}
  \item $\sigma_{x'}$ must lie between the principal stresses
        $\sigma_2 = \sigma_{\text{avg}} - R = -\SI{10.94}{ksi}$ and
        $\sigma_1 = \sigma_{\text{avg}} + R = \SI{6.94}{ksi}$. \checkmark
  \item $|\tau_{x'y'}|$ must be no larger than $R = \SI{8.94}{ksi}$. \checkmark
  \item The second principal direction is
        $\theta_{p2} = \ang{90} - |\theta_{p1}| = \ang{76.72}$, only $\ang{1.72}$ away from
        our cut at $\theta = \ang{75}$. So the cut face is \emph{almost} a principal plane,
        which is exactly why $\sigma_{x'} \approx \sigma_2$ and $\tau_{x'y'} \approx 0$.
        (On the circle, $X'$ sits just $\alpha = 2 \times \ang{1.72}$ away from $\sigma_2$.)
\end{itemize}

% ==================================================================
\section*{Practice: Problems 2 and 3}
% ==================================================================

\noindent
Use the same four steps to find $\sigma_{x'}$ and $\tau_{x'y'}$ on the inclined faces
shown. Start with Step~0: sketch the $x'$ (outward normal) and $y'$ directions on each cut.

\begin{center}
\begin{minipage}[t]{0.45\textwidth}
\centering
\textbf{Problem 2}\\[0.6em]
\begin{tikzpicture}
  \draw[thick, fill=gray!12] (0,0) rectangle (2,2);
  \draw[thick] (0.6,2) -- (2,0);
  \draw[thinarrow] (1.45,0) arc (180:125:0.55);
  \node at (1.05,0.22) {\small$\ang{55}$};
  % sigma_x = -6 ksi (compression), sigma_y = +12 ksi (tension), no shear
  \draw[force] (-0.8,1) -- (-0.05,1);
  \draw[force] (2.8,1) node[above] {$\SI{6}{ksi}$} -- (2.05,1);
  \draw[force] (1,2.05) -- (1,2.8) node[right] {$\SI{12}{ksi}$};
  \draw[force] (1,-0.05) -- (1,-0.8);
\end{tikzpicture}\\[0.6em]
\textbf{Answer:} $\sigma_{x'} = -\SI{0.078}{ksi}$,\; $\tau_{x'y'} = \SI{8.457}{ksi}$
\end{minipage}%
\hfill
\begin{minipage}[t]{0.5\textwidth}
\centering
\textbf{Problem 3}\\[0.6em]
\begin{tikzpicture}
  \draw[thick, fill=gray!12] (0,0) rectangle (2,2);
  \draw[thick] (0,2) -- (0.728,0);
  \draw[thinarrow] (0.55,2) arc (0:-70:0.55);
  \node at (0.85,1.55) {\small$\ang{70}$};
  % sigma_x = 120 MPa, sigma_y = 60 MPa, tau_xy = +45 MPa
  \draw[force] (2.05,1) -- (2.8,1) node[below right] {$\SI{120}{MPa}$};
  \draw[force] (-0.05,1) -- (-0.8,1);
  \draw[force] (1,2.05) -- (1,2.8) node[right] {$\SI{60}{MPa}$};
  \draw[force] (1,-0.05) -- (1,-0.8);
  \draw[force] (2.25,0.4) -- (2.25,1.6);
  \draw[force] (-0.25,1.6) -- (-0.25,0.4);
  \draw[force] (0.4,2.25) -- (1.6,2.25) node[right] {$\SI{45}{MPa}$};
  \draw[force] (1.6,-0.25) -- (0.4,-0.25);
\end{tikzpicture}\\[0.6em]
\textbf{Answer:} $\sigma_{x'} = \SI{141.90}{MPa}$,\; $\tau_{x'y'} = \SI{15.18}{MPa}$
\end{minipage}
\end{center}

\end{document}
